%=====================================================================
\section{Discusión y limitaciones}

\textbf{El costo de la interpretabilidad es medible y moderado.} Con 3 a 6 reglas legibles y unos 27 parámetros, el modelo difuso conserva el \cifCuno{} del AUC-PR de un \emph{gradient boosting} con miles de nodos, y empata en rango medio con la regresión logística. La hipótesis de que la interpretabilidad intrínseca daría además robustez a la deriva no se sostiene: el difuso cae más que las referencias.

Una explicación consistente con los errores es que las reglas apuestan casi todo al tamaño y a la difusión. Cuando la deriva cambia la mezcla de commits (más correcciones pequeñas, más revisión), un modelo con pocas reglas no tiene con qué compensar.

\textbf{Amenazas a la validez.}
\begin{itemize}
\item \textbf{La prueba no estuvo intacta.} En la fase de factibilidad inicial se exploró 2026 con modelos simples. Ninguna decisión se basó en eso, pero la prueba no es virgen. La iteración 2 se decidió sin volver a evaluarla.
\item \textbf{La etiqueta.} \texttt{Fixes:} es fiable en torno al 81\,\% \cite{Lyu2024}. La censura de 90 días ignora los errores que se descubren más tarde \cite{Falessi2022}.
\item \textbf{El subgrupo \texttt{Assisted-by}} es una cota inferior: el 90\,\% de los commits asistidos no la declaraba \cite{KernelRemove}.
\item \textbf{Validez externa.} Se usó un solo proyecto.
\item \textbf{La búsqueda.} El mejor punto quedó en el extremo del espacio (población y generaciones máximas), así que podría haber margen.
\item \textbf{La iteración 2} comparó la semilla 11 contra un umbral calculado con 5 semillas. Aun con la semilla 11, ningún candidato superó a la iteración 1, así que la decisión no depende de esa asimetría.
\item \textbf{McNemar} compara aciertos con un umbral fijo y, con este desbalance, castiga a quien marca más. Por eso se complementa con el AUC-PR y sus intervalos.
\end{itemize}

%=====================================================================
\section{Conclusiones}

\begin{enumerate}
\item El clasificador difuso-evolutivo alcanzó la meta de desempeño de C1 al lograr en 2026 \cifTestFE{} de AUC-PR (\cifCuno{} del \emph{gradient boosting}, superando la meta del 90\,\%), aunque cayó \cifCaidaFE{} desde validación frente a \cifCaidaGB{} de la referencia.
\item Aun así, supera con amplio margen al azar (IC \cifTestFEic{} frente a una prevalencia de \cifPrevPrueba{}). Lo hace con \cifNReglasRef{} reglas que un mantenedor puede leer, frente a miles de parámetros.
\item La evolución es lo que hace útil al modelo: \cifAporteEvolPrueba{} más de AUC-PR que la variante sin evolución y, a la vez, de 15 a 3--6 reglas.
\item La debilidad práctica es el umbral. Con pocas reglas el puntaje es escalonado, y un corte fijado en 2025 no se transfiere a 2026. Recalibrarlo con datos recientes es el siguiente paso natural, junto con aptitudes que premien la estabilidad temporal.
\item Ninguno de los modelos predice bien los commits declarados como asistidos por IA (AUC-PR \cifAssistedFE{}--\cifAssistedLR{}), una pregunta abierta para cuando haya una etiqueta fiable.
\end{enumerate}
